% typography.tex — improvement #79
%
% The three vendored families, and the body text settings that depend on them.
%
% Before this item the build passed `--variable=fontsize:10pt` and nothing else, so
% tectonic fell back to Latin Modern — a 1970s Computer Modern revival that sets thin and
% grey at 10pt on uncoated stock. Every face below is loaded from assets/fonts/ by path,
% never by system name, so the build is reproducible on a machine with no fonts installed.

\usepackage{fontspec}

% Path and Extension are set **per family**, never with \defaultfontfeatures. A global
% Path sends every later font lookup into assets/fonts/ — including the template's own
% \setmathfont{latinmodern-math}, which then fails to resolve and takes the build down at
% \begin{document} with an error that names a font nothing here asked for.
\newcommand{\bookfontpath}{assets/fonts/}

% --- Body: Source Serif 4 --------------------------------------------------
%
% The default optical size, not SmText or Caption: those are cut for 6–8.5pt, and the book
% sets its body at 10pt. Semibold is here for the rare bold run inside body text, where
% Bold is heavier than the page wants.

\setmainfont{SourceSerif4}[
  Path = \bookfontpath,
  Extension = .otf,
  Ligatures = TeX,
  UprightFont = *-Regular,
  ItalicFont = *-It,
  BoldFont = *-Bold,
  BoldItalicFont = *-BoldIt,
  FontFace = {sb}{n}{Font=SourceSerif4-Semibold},
  Numbers = OldStyle,
]

% Old-style figures in running text; lining figures where numbers have to align — a table
% column, a version stamp, a folio. fontspec already provides `\liningnums{…}` for that,
% so structure.tex and blocks.tex use it rather than a second command that would shadow it.

% --- Display: Source Sans 3 ------------------------------------------------
%
% Carries every small structural label in the book: eyebrows, running heads, table
% headers, code labels, callout titles. It has to stay legible at 8pt, which is where a
% humanist sans with a tall x-height earns its place over the serif.

\setsansfont{SourceSans3}[
  Path = \bookfontpath,
  Extension = .otf,
  Ligatures = TeX,
  UprightFont = *-Regular,
  ItalicFont = *-It,
  BoldFont = *-Bold,
  BoldItalicFont = *-BoldIt,
  FontFace = {sb}{n}{Font=SourceSans3-Semibold},
  FontFace = {sb}{it}{Font=SourceSans3-SemiboldIt},
  Numbers = Lining,
]

% Semibold is the weight the design actually calls for on labels; Bold is reserved for
% chapter and part titles. Naming it once here keeps that distinction honest.
\newcommand{\labelfont}{\sffamily\fontseries{sb}\selectfont}

% --- Code: Source Code Pro -------------------------------------------------
%
% Same family metrics as the other two, so a code span inside a sentence does not change
% the line's colour. Scaled to match the serif's x-height rather than its point size.

\setmonofont{SourceCodePro}[
  Path = \bookfontpath,
  Extension = .otf,
  Ligatures = TeX,
  UprightFont = *-Regular,
  ItalicFont = *-It,
  BoldFont = *-Bold,
  BoldItalicFont = *-BoldIt,
  Scale = MatchLowercase,
  Numbers = Lining,
]

% --- Body text settings ----------------------------------------------------

\usepackage{microtype}

% The class is asked for 10pt by build-book.sh, which gives a 12pt baseline. The design
% wants \tokBodyLeading, so the ratio between the two is what gets set — a \fontsize here
% would be undone by the first \normalsize.
\linespread{\tokLeadRatio}\selectfont

% Paragraphs are separated by space rather than indented — the book is full of lists,
% tables and code, and an indent after every one of them reads as an error. The value is a
% token because it is worth roughly ninety pages across the book.
\setlength{\parindent}{0pt}
\setlength{\parskip}{\tokParSkip}

% Ragged right, not justified. A 72–90 character measure with this much inline code in it
% cannot be justified without rivers, and TeX's alternative is an overfull box per
% paragraph. `ragged2e` keeps hyphenation on, which plain \raggedright turns off.
\usepackage[newcommands]{ragged2e}
\RaggedRight

% Long code lines and wide table cells would otherwise run into the margin.
\setlength{\emergencystretch}{3em}
